\chapter{Conclusion and Future Work}
\label{ch:conclusion}

\section{Summary of Contributions}
\label{sec:conclusion_summary}

This thesis presented a privacy-preserving, federated anomaly detection
framework for distributed microservice architectures, integrating six core
components: (1)~a federated LSTM-Autoencoder achieving F1~$=$~0.839 and
AUC-ROC~$=$~0.950 on real-world telemetry; (2)~differential privacy with
empirically characterized gradient clipping and noise trade-offs;
(3)~causal-graph-based root cause analysis with 89\% Top-3 accuracy;
(4)~Q-learning adaptive thresholding with convergence analysis;
(5)~an adaptive compression pipeline achieving up to 61\% payload reduction;
and (6)~fault-tolerant edge design with two node profiles.

\section{Conclusion}
\label{sec:conclusion_conclusion}

The experimental evaluation yields the following key results:

\begin{itemize}
    \item \textbf{High accuracy on real-world data.} The federated LSTM-AE
    achieves F1~$=$~0.839, outperforming the centralized baseline (0.496).
    Federated aggregation of distributed representations provides stronger
    generalization than single-node centralization. The model operates on
    ten selected telemetry features spanning CPU, memory, network, and
    application health---a feature set designed to capture complementary
    signals across infrastructure and application layers.

    \item \textbf{Sharp DP phase transition.} AUC-ROC remains near-baseline
    (0.970--0.971) for $\sigma \leq 0.0005$ ($\varepsilon \geq 10$). A sharp
    phase transition occurs at $\sigma{=}0.001$ ($\varepsilon \approx 1$),
    where F1 drops from 0.859 to 0.235 with default clipping ($C{=}0.5$).

    \item \textbf{Clipping norm as a critical DP lever.} The clipping
    norm has a dramatic effect: at $C{=}0.01$ with $\sigma{=}0.001$,
    the model achieves F1~$=$~0.868 and AUC-ROC~$=$~0.976---matching
    the no-DP baseline. Since noise is scaled by $\sigma \cdot C$, tighter
    clipping effectively neutralizes DP noise while preserving formal privacy
    guarantees. This resolves the traditionally assumed privacy--utility
    tension for this class of anomaly detection models.

    \item \textbf{Q-learning cold-start.} The Q-learning agent maintains a
    conservative threshold during the 24-hour evaluation, producing zero false
    positives but also zero true positives. The reward dominance from true
    negatives prevents threshold adaptation under extreme class imbalance (2\%
    anomaly rate), identifying a clear limitation addressable through
    asymmetric reward tuning or deep Q-networks.

    \item \textbf{Efficient communication.} Binary serialization produces
    payloads $5.4\times$ smaller than text formats. The full FL pipeline
    accounts for less than 3\% of round time, with local training dominating
    at over 97\%.

    \item \textbf{Actionable RCA.} The RCA engine places the true root cause
    in its top-3 ranking for 89\% of fault-injection events, with MRR of
    0.653 outperforming the random baseline (0.601).

    \item \textbf{Linear scalability.} FedAvg aggregation scales linearly
    from 5.1\,ms at 5~nodes to 990.1\,ms at 1\,000~nodes with bounded memory.
\end{itemize}

\section{Limitations}
\label{sec:conclusion_limitations}

\begin{itemize}
    \item \textbf{Q-learning cold-start.} The tabular Q-learning agent
    requires extended training horizons (${>}$24h) and asymmetric reward
    functions before its threshold policy stabilizes, during which alert
    quality may be suboptimal.

    \item \textbf{IID data assumption.} FedAvg assumes approximately IID
    local datasets \cite{li2020federated}. Production deployments with
    heterogeneous workloads may require alternative aggregation strategies
    such as FedProx or scaffold methods.

    \item \textbf{RCA on deep chains.} Top-1 accuracy of 42\% indicates
    room for improvement, particularly for services at the bottom of deep
    dependency chains where multiple propagation paths exist.

    \item \textbf{Clipping--privacy tension.} While tight clipping
    ($C{=}0.01$) recovers accuracy, it reduces the effective noise magnitude
    to $\sigma \cdot C = 10^{-5}$, weakening the formal privacy guarantee.
    The $(\varepsilon, \delta)$-DP bound improves numerically, but the
    practical protection against reconstruction attacks requires further
    empirical validation.

    \item \textbf{Single coordinator.} Despite crash-recovery mechanisms,
    the coordinator remains a potential bottleneck under extreme node counts.
\end{itemize}

\section{Future Work}
\label{sec:conclusion_future}

\textbf{Transfer Learning.} Pre-training on large telemetry corpora (e.g.,
Alibaba Cluster Trace) followed by few-shot fine-tuning would reduce FL
convergence time for newly onboarded services and mitigate the cold-start
problem.

\textbf{Hierarchical Federated Learning.} A two-tier architecture with
cluster-level sub-coordinators would extend scalability beyond 1\,000~nodes
\cite{kairouz2021advances}.

\textbf{Enhanced Privacy Mechanisms.} Secure Multi-Party Computation,
homomorphic encryption, and Local Differential Privacy would strengthen the
threat model beyond honest-but-curious coordinators. Combining these with the
tight-clipping strategy identified in this work could provide strong formal
guarantees without accuracy loss.

\textbf{Improved RCA.} Multi-modal causal graphs incorporating log events
and deployment metadata, temporal edge annotations, and LLM-assisted incident
summaries would improve accuracy beyond the current 42\% Top-1.

\textbf{Deep Q-Network.} Replacing tabular Q-learning with DQN would enable
continuous state representation, richer features (rolling statistics,
time-of-day encoding), and faster convergence via experience replay. This
directly addresses the cold-start limitation by enabling batch learning from
stored transition tuples.

\textbf{Multi-Modal Anomaly Detection.} Extending the LSTM-AE to jointly
model metrics, structured logs, and traces through a multi-modal encoder
could improve detection of complex anomalies that manifest across multiple
telemetry types simultaneously.
